\chapter{The Web Application}
\label{ch:web}

The web client is \voname{}'s primary collaboration surface: a Slack-style
single-page application through which a member manages workspaces, chats in
channels and direct messages, opens threads, joins room calls with live
voice/video, and receives notifications. This chapter describes its architecture,
its authentication and API strategy, its three-connection real-time layer, its
Agora-based calling, and a feature-by-feature tour of the application.

\section{Overview and Technology Stack}

The client is a React~18 single-page application created with Create React App.
It uses \textbf{MUI~7} and \textbf{Tailwind CSS} for its interface,
\textbf{axios} and the browser \code{fetch} API for REST, \textbf{\code{@stomp/
stompjs}} for real-time messaging, and \textbf{\code{agora-rtc-sdk-ng}} for
voice and video. State management is deliberately minimal: there is no Redux or
global context. Component-local state is lifted into the top-level chat page and
passed down by props, session and theme live in \code{localStorage}, and the
notification socket is a module singleton. This keeps the codebase small and the
data flow easy to follow.

\begin{notebox}
The web client is the Slack-style collaboration app, \emph{not} the 2D office;
the animated world (Chapter~\ref{ch:skyoffice}) is a separate client. Within the
web app, the ``office'' metaphor is realised through workspaces and through
\code{ROOM}-type channels that host Agora audio/video calls.
\end{notebox}

\section{Application Structure and Routing}

The application boots by mounting the router inside a strict-mode root, with an
early theme bootstrap that reads the saved (or system) colour scheme and toggles
a dark-mode class on the document body --- the mechanism the whole app uses for
theming. Routes divide into public and protected (Table~\ref{tab:web-routes}); a
lightweight \code{PrivateRoute} guard redirects to login when no token is present.

\begin{table}[htbp]
\centering
\caption{Web client routes.}
\label{tab:web-routes}
\small
\begin{tabularx}{\textwidth}{L{3.4cm} L{2.0cm} X}
\toprule
\textbf{Path} & \textbf{Access} & \textbf{Purpose} \\
\midrule
\route{/login}, \route{/signup} & Public & Authentication \\
\route{/verify-email} & Public & OTP e-mail verification \\
\route{/chat} & Protected & The main application (workspaces, channels, rooms, DMs, threads, calls) \\
\route{/chat/join/:channelId} & Protected & Invite-link landing; joins a channel then redirects \\
\bottomrule
\end{tabularx}
\end{table}

\section{Authentication and the API Layer}

On login or registration the client stores the returned JWT in
\code{localStorage} and navigates into the application. It reads the current
user's numeric id by decoding the JWT payload on the client side --- that id is
sent as the \code{X-User-Id} header on nearly every subsequent request. Two REST
conventions coexist: an \textbf{axios instance} (used for auth and profile) whose
interceptors attach the bearer token and, on any \code{401}, clear the session and
redirect to login; and \textbf{raw \code{fetch}} calls (used for chat, workspace,
rooms, and notifications) that set the identity headers by hand.

During development, a proxy configuration doubles as the service map, forwarding
each path prefix to the corresponding microservice port and, importantly, wiring
the \emph{WebSocket upgrade} for the chat and room sockets
(Table~\ref{tab:web-proxy}). In production these resolve to the gateway.

\begin{table}[htbp]
\centering
\caption{Development proxy: the frontend's view of the backend.}
\label{tab:web-proxy}
\small
\begin{tabularx}{\textwidth}{L{3.2cm} L{2.0cm} X}
\toprule
\textbf{Path prefix} & \textbf{Port} & \textbf{Service} \\
\midrule
\route{/api/auth}, \route{/api/users} & 8091 & User \\
\route{/api/chat} (+WS) & 8084 & Chat \\
\route{/api/rooms}, \route{/ws/rooms} & 8086 & Room \\
\route{/api/workspace} & 8087 & Workspace \\
\route{/api/notifications} & 8082 & Notifications \\
\bottomrule
\end{tabularx}
\end{table}

\section{The Real-time Layer: Three STOMP Connections}

The web client maintains \emph{three} independent STOMP connections, one to each
real-time service, all following the ticket pattern of
Section~\ref{sec:realtime-arch}: before every (re)connection a fresh single-use
ticket is fetched, and a five-second reconnect delay provides automatic recovery.

\begin{description}
  \item[Chat (8084).] The core connection. The message list for the active
  channel subscribes to \route{/topic/channel/\{id\}} and handles the
  \code{NEW\_MESSAGE}, \code{EDIT\_MESSAGE}, and \code{DELETE\_MESSAGE} actions,
  merging live events with fetched history by id. The composer publishes to
  \route{/app/chat/send}. Threads use the parallel \route{/topic/thread/\{id\}}.
  \item[Rooms (8086).] When the active channel is a room, the client subscribes
  to \route{/topic/room/\{roomId\}} for call presence
  (\code{PARTICIPANT\_JOINED}/\code{LEFT}, \code{STATE\_CHANGED}) and publishes its
  own join/leave to \route{/app/room/state}, so the ``N in call'' indicator stays
  live for everyone.
  \item[Notifications (8082).] Implemented as a reference-counted module
  singleton so the whole app shares one socket; it subscribes to
  \route{/user/queue/notifications} and raises the unread badge on each
  \code{NEW\_NOTIFICATION}.
\end{description}

Because every service uses the same \code{\{action, payload\}} envelope, each of
these subscriptions is handled by the same simple pattern --- switch on the
action, apply the payload --- which is a direct payoff of the architectural
uniformity described in Chapter~\ref{ch:architecture}.

\section{Voice and Video with Agora}

Real-time media in the web client is room-scoped and handled by a single reusable
hook that encapsulates the entire Agora engine. The flow to join a call
(Figure~\ref{fig:web-call}) is:

\begin{enumerate}
  \item The user opens a \code{ROOM}-type channel and clicks \emph{Join Call}.
  \item The client resolves the real room id and calls
  \route{POST /api/rooms/\{id\}/join}, which returns the Agora channel name and a
  freshly minted access token --- the room service is the token server
  (Section~\ref{sec:be-room}).
  \item The hook joins the Agora channel with that token, creates and publishes
  the local microphone and camera tracks, and subscribes to remote participants
  as they publish, rendering each into a video tile.
  \item In parallel, the client broadcasts its presence over the rooms STOMP
  connection so other members see the participant list update.
\end{enumerate}

\begin{figure}[htbp]
\centering
\begin{tikzpicture}[font=\footnotesize,
  b/.style={draw=voPrimary,fill=voPrimary!8,rounded corners=2pt,minimum height=0.7cm,
            text width=2.1cm,align=center},
  a/.style={-{Stealth[]},thick,voDark},note/.style={font=\scriptsize\itshape,voGray}]
\node[b] (c) at (0,0) {Web client};
\node[b] (r) at (4.4,0) {Room Svc};
\node[b] (ag) at (9.2,0) {Agora SFU};
% client <-> room service: two separated arcs
\draw[a] (c.20) to[out=22,in=158] node[above,note]{POST /rooms/id/join} (r.160);
\draw[a] (r.200) to[out=-158,in=-22] node[below,note]{channel + token} (c.-20);
% client <-> Agora SFU: two arcs that bow clear of the Room Svc box
\draw[a] (c.-42) to[out=-38,in=214] node[below,note,pos=0.5]{publish mic + camera} (ag.218);
\draw[a] (ag.142) to[out=142,in=38] node[above,note,pos=0.5]{remote tracks} (c.42);
\end{tikzpicture}
\caption{Joining a room call: the room service mints the Agora token, media flows
client-to-SFU.}
\label{fig:web-call}
\end{figure}

Toggling the microphone or camera flips the enabled state of the local tracks;
leaving closes the tracks, leaves the channel, and broadcasts departure. Crucially
the media never traverses \voname{}'s own servers --- the backend's only role is
to authorise the join by issuing a scoped token.

\section{Feature Tour}

The protected \route{/chat} route is the composition root and holds nearly all
application state. Its responsive four-column layout --- workspace rail, channel
sidebar, chat area, and (when open) thread panel and member list --- adapts down
to a drawer-based mobile layout.

\begin{description}
  \item[Workspaces.] A vertical rail lists the user's workspaces (loaded from the
  workspace service); if the user has none, a default workspace is created
  automatically. A dialog creates new workspaces with a name, auto-generated slug,
  description, and timezone.
  \item[Channels, rooms, and DMs.] The sidebar lists group channels, room
  channels, and direct messages for the active workspace, auto-provisioning a set
  of starter channels for a brand-new workspace. Selecting a channel loads its
  history and subscribes to its live topic; selecting a room additionally enables
  the call interface. Direct messages are opened by target user id and guaranteed
  unique per pair by the backend.
  \item[Messaging.] The chat area renders the merged history-plus-live message
  stream with date dividers, resolves sender names through a client-built user
  map, and supports inline editing and deletion of one's own messages, replies,
  mentions, an emoji picker, and a client-side message search. A ``reply in
  thread'' action opens the slide-in thread panel, which creates or loads a thread
  and subscribes to its topic.
  \item[Rooms and calls.] For room channels, a call bar and a video grid layer the
  Agora experience described above on top of the same message timeline, with a
  live participant count driven by the rooms socket.
  \item[Notifications.] A bell in the sidebar footer shows the unread count, opens
  a list fed by REST plus the live notifications socket, and supports mark-read
  and delete.
  \item[Settings and invites.] A settings dialog handles avatar upload and
  password change; channels can be shared by a copyable invite link that lands on
  the \route{/chat/join/:channelId} route and auto-selects the joined channel.
\end{description}

\section{Design Observations}

Several deliberate choices are worth noting for the evaluation in
Chapter~\ref{ch:testing}. The ticket-based WebSocket authentication, refreshed on
every reconnect, is a clean and robust solution to the browser's inability to set
socket headers. The single \code{\{action, payload\}} dispatch pattern keeps the
real-time code uniform across three sockets. On the other hand, the client is
honestly at \emph{prototype} maturity in places: some creation flows use native
browser prompts, presence beyond ``in this call'' is not yet backed by a server
roster, the typing indicator is sent but not yet displayed, and message
attachments are rendered but not yet uploadable. These are precisely the rough
edges catalogued as future work.
